\documentclass[11pt]{article}

% parameters are fairly self-explanatory here, thankfully
\usepackage[width=7in,height=9.5in,top=0.75in,centering]{geometry}

% for math-ing
\usepackage{amsmath}
\usepackage{amssymb}


%%%%%% tables %%%%%%%%

\usepackage{array}
\renewcommand{\arraystretch}{1.5}
\setlength{\tabcolsep}{8pt}

%%%%%% plotting %%%%%%

\usepackage[dvipsnames]{xcolor}
\usepackage{tikz}
\usepackage{pgfplots}

\pgfplotsset{every axis/.append style={
                    axis x line=middle,    % put the x axis in the middle
                    axis y line=middle,    % put the y axis in the middle
                    axis line style={<->,color=black}, % arrows on the axis
                    xlabel={$x$},          % default put x on x-axis
                    ylabel={$y$},          % default put y on y-axis
            }}
            
 \pgfplotsset{compat=1.14} % sometimes, everything falls apart without this
 
 %\usepgfplotslibrary{fillbetween} %for shading between curves

%%%%%% commands %%%%%%

% exam heading; course name is first argument, exam information is second argument
\newcommand{\Heading}[2]{
    \fbox{
        {\Large\textbf{#1}}
        }
    \hfill
    \fbox{
        \begin{minipage}{0.2\textwidth}
            \begin{center}
            \textbf{#2}
            \end{center}
        \end{minipage}
    }
    \par\vspace{0.1in}}

% for being lazy while beautifying integrals
\newcommand{\dint}{\displaystyle\int}

% for being lazy while beautifying limits
\newcommand{\dlim}{\displaystyle\lim}

%%%%%% stuff %%%%%%

\usepackage{fancyhdr}

%\setcounter{page}{0} % first page is page 0

\parindent0pt % no indent

\fboxsep=0.25cm % little space in fbox border

%%%%%% BEGIN! %%%%%%%

\begin{document}

\pagestyle{fancy}

% record goals here to create sense of transparency

\fancyhead[l]{%
  \ifcase\value{page}
  % empty test for page = 0
  \or {}% 
  \or
  \or {\bf\color{ProcessBlue} F1: Lines \& Slope}% page=1
  \or {\bf\color{RoyalBlue} L1: Limits and Continuity}% page = 2
  \or {\bf\color{RoyalBlue} L2: Evaluating Limits Using Continuity and Forms}% page = 3
  \or {\bf\color{ForestGreen} D1: Meaning of a Derivative}% page = 4
  \or {\bf\color{ProcessBlue} F2: Exponential and Logarithmic Functions}% page = 5
  \or {\bf\color{ForestGreen} D2: Compute Derivatives}% page = 6
  \or {\bf\color{ForestGreen} D3: Derivatives of Basic Functions}% page = 7
  \or {\bf\color{RedOrange} A1: Derivatives and Graphing}% page = 8
  \or
  \else
  % Empty chead here!
  \fi
}

\thispagestyle{empty} % don't display that this is page 0

\Heading{MATH 1190/1210 Survey of Calculus I}{Quiz \\Learning Target P2\\Fall 2025}

\vspace{0.1in}

\textbf{Name} \hrulefill \, \begin{minipage}{0.16\textwidth}\begin{center}\textbf{UVA ID\\ (e.g. hdj4nd)} \end{center}\end{minipage}\hrulefill\\




%%%%%%%%%%%%%%%%%%%%%%%%%%%%%%%%%%%%%%%%%%%%%%%%%%%%%
%%%%%%%%%%% BEGIN PROBLEM A1 %%%%%%%%%%%%%%%%%%%%%%%%
%%%%%%%%%%%%%%%%%%%%%%%%%%%%%%%%%%%%%%%%%%%%%%%%%%%%%
This problem is based on Lesson 22. Questions about Mean and Variance would not appear on the KC4, but will appear on the final!

Let $f(x)=x  + \dfrac{1}{x}$. You may take as given that $f'(x) = 1-\dfrac{1}{x^2} = \dfrac{x^2-1}{x^2}$ and $f''(x)=\dfrac{2}{x^3}$.

\begin{enumerate}
    \item Find all critical numbers of $f$. Find intervals where $f(x)$ is increasing. 
    \vfill
    \item Find $x$-coordinates of all inflection points.
    \vfill
\end{enumerate}

%%%%%%%%%%%%%%%%%%%%%%%%%%%%%%%%%%%%%%%%%%%%%%%%%%%%%
%%%%%%%%%%% END PROBLEM A1 %%%%%%%%%%%%%%%%%%%%%%%%%%
%%%%%%%%%%%%%%%%%%%%%%%%%%%%%%%%%%%%%%%%%%%%%%%%%%%%%


\end{document}